\documentclass[12pt,oneside]{book}

% ============================================================
% METADATOS
% ============================================================
\newcommand{\universidad}{Universidad de Buenos Aires (UBA)}
\newcommand{\materia}{Derecho Procesal Civil}
\newcommand{\titulo}{Unidad 25 --- Medidas Cautelares en Particular}
\newcommand{\autor}{Isaac Edgar Camacho Ocampo}
\newcommand{\anio}{2024}

% ============================================================
% PAQUETES
% ============================================================
\usepackage[utf8]{inputenc}
\usepackage[T1]{fontenc}
\usepackage[spanish,es-nodecimaldot]{babel}

\usepackage[
    top=2.5cm,
    bottom=2.5cm,
    left=3cm,
    right=2.5cm,
    headheight=15pt
]{geometry}

\usepackage{amsmath}
\usepackage{amssymb}
\usepackage{mathtools}

\usepackage{graphicx}
\usepackage{float}
\usepackage{caption}
\usepackage{subcaption}

\usepackage{booktabs}
\usepackage{array}
\usepackage{longtable}
\usepackage{tabularx}

\usepackage{enumitem}

\usepackage{xcolor}
\usepackage[most]{tcolorbox}

\usepackage{fancyhdr}

\usepackage{ragged2e}

\usepackage[
    colorlinks=true,
    linkcolor=blue!50!black,
    citecolor=green!40!black,
    urlcolor=blue!60!black,
    pdftitle={\titulo},
    pdfauthor={\autor},
    pdfsubject={Apuntes universitarios},
    bookmarks=true
]{hyperref}

% ============================================================
% RUTAS DE IMÁGENES
% ============================================================
\graphicspath{
    {./img/}
    {./graficos/}
    {./figuras/}
}

% ============================================================
% CONFIGURACIÓN
% ============================================================
\setcounter{tocdepth}{2}

\setlength{\parindent}{0pt}
\setlength{\parskip}{0.5em}

\setlist[itemize]{
    topsep=0.3em,
    itemsep=0.2em
}

\setlist[enumerate]{
    topsep=0.3em,
    itemsep=0.2em
}

\clubpenalty=10000
\widowpenalty=10000

% ============================================================
% COMANDOS MATEMÁTICOS
% ============================================================
\newcommand{\abs}[1]{\left|#1\right|}
\newcommand{\norm}[1]{\left\lVert#1\right\rVert}
\newcommand{\vect}[1]{\mathbf{#1}}
\newcommand{\deriv}[2]{\frac{d#1}{d#2}}
\newcommand{\pderiv}[2]{\frac{\partial#1}{\partial#2}}

% ============================================================
% ENCABEZADOS Y PIES
% ============================================================
\pagestyle{fancy}
\fancyhf{}

\fancyhead[L]{\nouppercase{\leftmark}}
\fancyhead[R]{\materia}

\fancyfoot[C]{\thepage}

\renewcommand{\headrulewidth}{0.4pt}
\renewcommand{\footrulewidth}{0pt}

\fancypagestyle{plain}{
    \fancyhf{}
    \fancyfoot[C]{\thepage}
    \renewcommand{\headrulewidth}{0pt}
}

% ============================================================
% CAJAS ACADÉMICAS
% ============================================================
\newtcolorbox{definition}[1]{
    enhanced,
    breakable,
    title={Definición: #1},
    fonttitle=\bfseries,
    colback=blue!3,
    colframe=blue!50!black,
    boxrule=0.8pt,
    arc=2mm
}

\newtcolorbox{important}{
    enhanced,
    breakable,
    title={Importante},
    fonttitle=\bfseries,
    colback=yellow!8,
    colframe=orange!70!black,
    boxrule=0.8pt,
    arc=2mm
}

\newtcolorbox{example}[1]{
    enhanced,
    breakable,
    title={Ejemplo: #1},
    fonttitle=\bfseries,
    colback=green!4,
    colframe=green!40!black,
    boxrule=0.8pt,
    arc=2mm
}

\newtcolorbox{formula}[1]{
    enhanced,
    breakable,
    title={Fórmula / Regla: #1},
    fonttitle=\bfseries,
    colback=purple!4,
    colframe=purple!50!black,
    boxrule=0.8pt,
    arc=2mm
}

\newtcolorbox{exercise}[1]{
    enhanced,
    breakable,
    title={Ejercicio: #1},
    fonttitle=\bfseries,
    colback=gray!5,
    colframe=gray!60!black,
    boxrule=0.8pt,
    arc=2mm
}

\newtcolorbox{note}{
    enhanced,
    breakable,
    title={Nota},
    fonttitle=\bfseries,
    colback=white,
    colframe=gray!50,
    boxrule=0.6pt,
    arc=2mm
}

% ============================================================
% DOCUMENTO
% ============================================================
\begin{document}

% ------------------------------------------------------------
% PORTADA
% ------------------------------------------------------------
\begin{titlepage}

\thispagestyle{empty}

\begin{center}

\vspace*{1cm}

{\large \textsc{\universidad}}

\vspace{3cm}

{\Huge \textbf{\materia}}

\vspace{1cm}

{\LARGE \titulo}

\vspace{0.8cm}

\textit{Comprender los conceptos es el primer paso para convertir el estudio en conocimiento.}

\vspace{1.2cm}

\rule{0.70\textwidth}{0.5pt}

\vspace{0.5cm}

{\large Apuntes de estudio}

\vspace{0.5cm}

\rule{0.70\textwidth}{0.5pt}

\vfill

{\large \autor}

\vspace{0.3cm}

{\large \anio}

\end{center}

\vfill

\begin{flushleft}
\textit{Este es un modesto aporte a los estudiantes de ``Derecho Procesal Civil'' no pretende sustituir el material oficial de la materia.}
\end{flushleft}

\end{titlepage}

% ------------------------------------------------------------
% ÍNDICE
% ------------------------------------------------------------
\tableofcontents

% ============================================================
% CONTENIDO
% ============================================================

\chapter{Introducción a las medidas cautelares en particular}

Esta unidad desarrolla las medidas cautelares en particular que forman parte del núcleo del proceso civil. Así como existen distintos instrumentos para tratar diferentes situaciones, el proceso civil cuenta con diversas herramientas para garantizar la eficacia de la sentencia: desde el embargo ---la más usada, análoga a una garantía sobre un bien--- hasta las medidas de familia, orientadas a proteger derechos fundamentales de los niños y los vínculos familiares.

\begin{important}
Cada una de las medidas cautelares en particular desarrolladas en esta unidad debe estudiarse en conexión directa con los presupuestos generales de las medidas cautelares (verosimilitud del derecho, peligro en la demora y contracautela), ya vistos previamente. Todas las medidas en particular están atravesadas por dichos presupuestos generales.
\end{important}

\chapter{El embargo}

\section{El embargo como medida cautelar por antonomasia}

\begin{definition}{Embargo}
La doctrina ha señalado generalmente que el embargo es la medida cautelar por antonomasia. Esto significa que es la más usada y la más habitual, y que consiste en la afectación e individualización de un bien ---mueble o inmueble, registrable o no--- al resultado de un proceso.
\end{definition}

El bien así individualizado a través del embargo queda vinculado directamente con el proceso, queda a disposición del juez y opera como garantía concreta para asegurar la efectividad del resultado del juicio. Ese bien queda directamente afectado ante la eventualidad de que resulte necesario proceder a la realización de bienes ---es decir, sacar a la venta bienes para pagar la deuda reclamada en el proceso---, siendo el bien fundamental a la hora de procederse a la subasta pública como modo de garantizar el pago de la deuda.

\section{Tipos de embargo}

Pueden reconocerse tres tipos de embargo en el ordenamiento procesal, según el estadio del proceso en el que se traben.

\begin{table}[H]
\centering
\caption{Tipos de embargo según el momento procesal}
\begin{tabularx}{\textwidth}{l l X}
\toprule
\textbf{Tipo} & \textbf{Momento} & \textbf{Normativa / Características} \\
\midrule
Embargo preventivo & Antes del proceso o durante su tramitación & Arts. 209 a 212 CPCCN. Eminentemente cautelar. Exige los tres recaudos: verosimilitud del derecho, peligro en la demora y contracautela suficiente. \\
\addlinespace
Embargo ejecutivo & Una vez admitido el título en el juicio ejecutivo & Se traba cuando el juez reconoce el título como admisible y libra la orden de intimación de pago. Individualiza el bien que será ejecutado si no se paga. \\
\addlinespace
Embargo ejecutorio & En el proceso de ejecución de sentencia & Se dicta cuando la sentencia está confirmada y se avanza en la concreción del pago. Es el paso previo a la subasta del bien. Art. 551 CPCCN (sentencia de trance y remate). \\
\bottomrule
\end{tabularx}
\end{table}

\subsection{El embargo preventivo}

El embargo preventivo es el eminentemente cautelar y puede trabarse aún antes del proceso o durante su secuela, siempre que se reúnan los tres recaudos clásicos de toda medida cautelar: la verosimilitud del derecho, el peligro en la demora y la contracautela adecuada o suficiente. Su regulación se encuentra en los siguientes artículos:

\begin{itemize}
    \item \textbf{Art. 209 CPCCN}: establece los supuestos de procedencia del embargo preventivo.
    \item \textbf{Art. 210 CPCCN}: amplía los supuestos y complementa el artículo anterior.
    \item \textbf{Art. 211 CPCCN}: regula el procedimiento para la traba del embargo.
    \item \textbf{Art. 212 CPCCN}: establece hipótesis específicas derivadas de la secuela del proceso.
\end{itemize}

\subsection{El embargo ejecutivo}

En el trámite del proceso de ejecución existe una primera etapa consistente en el análisis del título. No siempre el título en sí mismo trae aparejada la ejecución; en ocasiones debe prepararse la vía ejecutiva, y en esos casos procede el embargo preventivo, porque todavía no existe un título ejecutivo declarado admisible.

Una vez que el juez advierte que se trata de un título \textit{prima facie} hábil ---ya sea porque se preparó la vía ejecutiva o porque objetivamente reúne todos los recaudos propios del título--- ese título se considera admisible y se libra la orden de intimación de pago. A partir de ese momento procede el embargo ejecutivo, sea porque se transforma el embargo preventivo trabado anteriormente en ejecutivo, sea porque se traba directamente como embargo ejecutivo.

\subsection{El embargo ejecutorio}

El embargo ejecutorio es el que se traba en el marco del proceso de ejecución de sentencia. Se dicta cuando la sentencia ya ha sido confirmada por el superior y se avanza en el proceso de concreción del pago derivado de una sentencia condenatoria. Ante sentencia firme y ejecutoriada, el embargo que se traba ya no es preventivo ---porque el derecho está consolidado y reconocido--- ni ejecutivo ---porque no se traba en el trámite de ejecución de un título ejecutivo--- sino ejecutorio, propio del ámbito de la ejecución de sentencia.

\begin{formula}{Transformación automática en embargo ejecutorio (art. 551 CPCCN)}
A partir del dictado de la sentencia de trance y remate (art. 551 CPCCN), los embargos anteriores trabados a lo largo del proceso se transforman en embargos ejecutorios por efecto de la ley, muchas veces sin disposición concreta del tribunal. El embargo ejecutorio constituye el paso previo a la subasta del bien.
\end{formula}

\section{Efectos de la traba del embargo}

El embargo es una medida procesal, pero produce efectos regulados en los códigos sustanciales, con consecuencias tanto civiles como penales que se ponen en juego en caso de incumplimiento de las obligaciones del depositario de la cosa embargada.

\begin{important}
El depositario del bien embargado tiene la obligación de conservar la cosa, no disminuir su valor y ponerla a disposición del tribunal cuando sea requerido. El Código Penal tipifica delitos vinculados al incumplimiento de las obligaciones del depositario.
\end{important}

Los efectos fundamentales del embargo son los siguientes:

\begin{itemize}
    \item El embargo individualiza el bien que queda afectado a la garantía y vinculado al proceso.
    \item El embargo no produce la desapropiación del bien por parte del afectado. El afectado puede usar y gozar de la cosa, en tanto y en cuanto no la haga perder valor.
    \item Tratándose de bienes registrables, el embargo no produce la imposibilidad de enajenación del bien: el bien puede ser vendido.
\end{itemize}

\subsection{El principio \textit{nemo plus iuris} y el embargo}

La venta de un bien embargado se rige por el principio \textbf{nemo plus iuris}, según el cual nadie puede transmitir sobre una cosa un derecho mayor al que posee sobre ella. Quien posee un bien embargado lo posee en tanto se encuentra embargado, de modo que si lo transfiere, lo transmite embargado, y quien lo adquiere adquiere sobre la cosa el mismo derecho que tenía su anterior titular.

\subsection{El plenario \textit{Sertop}: doctrina judicial obligatoria}

\begin{important}
\textbf{Plenario ``Sertop''} (Cámara Nacional de Apelaciones en lo Civil): constituye doctrina judicial obligatoria en la Ciudad Autónoma de Buenos Aires. Establece que el adquirente de un bien embargado asume la totalidad de la deuda ---no solo el monto por el cual se trabó la medida--- incluyendo ampliaciones, intereses y costas del proceso.

En el fuero \textbf{comercial} este criterio no rige: allí la transmisión se limita al monto de la cautela.
\end{important}

Dada la incertidumbre sobre el criterio aplicable a este tipo de transmisiones, resulta difícil que existan personas dispuestas a adquirir bienes embargados. El embargo suele operar así como una traba a la libre circulación de los bienes: quienes los adquieren generalmente especulan pagando un precio significativamente menor al del mercado, o tienen algún tipo de vinculación con el sujeto embargado.

\section{Supuestos de procedencia del embargo preventivo}

La enumeración legal de los supuestos de procedencia del embargo preventivo (arts.\ 209 y 210 CPCCN) es meramente enunciativa y orientativa, no taxativa: el código establece presupuestos en los que procede el embargo, pero las situaciones fácticas no deben necesariamente encuadrarse en alguno de ellos.

\begin{itemize}
    \item \textbf{Inciso 1° (art. 209)}: deudor sin domicilio en la República. Se presume el peligro en la demora porque puede ser dificultoso encontrarle bienes para responder.
    \item \textbf{Supuesto 2}: reconocimiento de la deuda a través de instrumento público. Se presume reunida la verosimilitud y el peligro en la demora.
    \item \textbf{Supuesto 3}: instrumento privado cuyas firmas son abonadas por dos testigos (información sumaria). El aval de dos testigos que declaran haber presenciado la firma o conocer la rúbrica del deudor permite dictar la medida sin mayores recaudos.
    \item \textbf{Supuesto 4}: contratos bilaterales instrumentados públicamente o en instrumento privado con dos testigos (información sumaria), que avalan que la prestación a cargo del acreedor fue cumplida y resta cumplir la contraprestación.
\end{itemize}

\subsection{Hipótesis procesales del artículo 212 CPCCN}

El artículo 212 CPCCN describe tres hipótesis en las que la secuela del proceso determina la procedencia del embargo preventivo:

\begin{itemize}
    \item \textbf{Inciso 1° (remisión al art. 63)}: declaración de rebeldía del demandado. El hecho de que la parte sea declarada rebelde da motivo formal para la traba del embargo preventivo.
    \item \textbf{Inciso 2°}: confesión de la existencia de la deuda en audiencia confesional, o declaración de confeso ficto (por no comparecer a la audiencia habiendo sido citado). En ambos casos procede el embargo sin mayores recaudos. Cuando el demandado no contesta la demanda, por efecto del art. 356 se presume la veracidad de los hechos lícitos descriptos, lo que también habilita la traba.
    \item \textbf{Inciso 3°}: existencia de sentencia definitiva condenatoria aunque esté apelada. Si el juez de primera instancia condenó al pago, el embargo preventivo está plenamente configurado; la eventual revisión por la Cámara no impide la medida.
\end{itemize}

\section{Bienes inembargables}

\begin{important}
El artículo 219 CPCCN enumera los bienes inembargables:
\begin{itemize}
    \item \textbf{Inciso 1°}: el lecho cotidiano del deudor, su esposa e hijos; las ropas y bienes muebles de indispensable uso, y las herramientas u objetos que utiliza el deudor para el ejercicio de su profesión u oficio. Criterio judicial: si el deudor posee, por ejemplo, cinco televisores, puede reconocerse la inembargabilidad de uno, quedando los demás embargables.
    \item \textbf{Inciso 2°}: los sepulcros, salvo que la deuda sea por el saldo del precio de compra del sepulcro o por bienes adquiridos para su construcción.
    \item \textbf{Inciso 3°}: los bienes declarados inembargables por leyes especiales (jubilaciones, pensiones, salarios, etc.).
\end{itemize}
\end{important}

\subsection{Inembargabilidad del salario}

El trabajador tiene, en su salario, un porcentaje de inembargabilidad establecido por la Ley de Contrato de Trabajo, que dispone que solo puede embargarse hasta el 20\% de sus ingresos.

\begin{example}{Límite del 20\% y deuda alimentaria}
Si el salario del deudor es \$30.000, el 20\% equivale a \$6.000. Sin embargo, si el deudor tiene tres hijos menores y la madre reclama alimentos, \$6.000 puede resultar insuficiente. En ese supuesto, el juez puede excepcionalmente elevar el límite de embargabilidad al 40\% o 50\% del ingreso, dejando de lado el límite de la Ley de Contrato de Trabajo, siempre como modo de garantizar la subsistencia del trabajador.
\end{example}

Un criterio similar rige para jubilaciones y pensiones, que en principio son inembargables, salvo en caso de deudas alimentarias, supuesto en el cual el juez puede establecer el porcentaje embargable garantizando la subsistencia del jubilado o pensionado.

\section{Caducidad del embargo}

En materia de embargo rige, como en general en todas las medidas cautelares, lo establecido por el artículo 207 CPCCN respecto de la caducidad de la traba.

\begin{formula}{Caducidad del embargo (art. 207 CPCCN)}
\begin{itemize}
    \item \textbf{Plazo}: 10 días desde la traba de la medida para iniciar el proceso principal.
    \item En procesos sujetos a mediación previa, el inicio de la mediación interrumpe el plazo de caducidad (equivale al inicio del proceso).
    \item \textbf{Embargos registrales} (Registro del Automotor): caducan a los \textbf{cinco años}. Debe pedirse la reinscripción dentro de ese plazo para mantener subsistente la medida.
\end{itemize}
\end{formula}

\chapter{El secuestro}

\section{Concepto y ámbito de aplicación}

\begin{important}
El secuestro solo opera respecto de bienes \textbf{muebles}. No existe secuestro de bienes inmuebles. Puede recaer sobre bienes muebles registrables (automotores, motocicletas) y no registrables. Los animales (ganado, caballos de carrera) también pueden ser objeto de secuestro.
\end{important}

El secuestro está orientado, como medida cautelar, a desapoderar al deudor de la disposición material y jurídica sobre una cosa de su propiedad.

\section{Tipos de secuestro}

\subsection{Secuestro como medida cautelar propiamente dicha}

El secuestro es excepcionalmente menos habitual que el embargo, pero de acuerdo con las características del bien, del deudor o de la situación en que se encuentra el bien, en algunas hipótesis procede el secuestro (art. 221 CPCCN) como único modo de garantizar la intangibilidad de la cosa. Procede cuando la cosa puede ser deteriorada por el uso continuo del deudor, o cuando este ha dado muestras de falta de cuidado y ha permitido que el bien pierda valor.

\begin{example}{Procedencia del secuestro como cautelar}
\textbf{Caso 1}: la cosa es una obra de arte almacenada en un depósito con humedad, que no recibe los cuidados indispensables para su preservación. En ese caso, es preferible el secuestro al embargo para evitar el deterioro del bien.

\textbf{Caso 2}: se trata de aparatología para realizar estudios técnicos o científicos que requiere mantenimiento que el deudor no realiza, generando un deterioro progresivo.
\end{example}

\subsection{Secuestro como sanción}

El secuestro también puede operar como sanción derivada de la conducta del deudor o de quien incumple una obligación procesal. El código lo prevé en dos hipótesis:

\begin{itemize}
    \item \textbf{Art. 128 CPCCN --- secuestro por retención de expediente}: cuando un letrado o un auxiliar del tribunal (escribano, perito) retira un expediente en préstamo y no lo devuelve en el plazo fijado por el juez, se lo intima a devolverlo; si no lo hace en el plazo de la intimación, procede el secuestro del expediente por oficial de justicia, con allanamiento de morada si fuera necesario.
    \item \textbf{Art. 323 CPCCN, incisos 1° a 5° --- diligencias preliminares}: en las diligencias preliminares puede pedirse a la parte demandada que exhiba un objeto, un inmueble, un testamento o los títulos de una cosa vendida. Si se incumple la intimación de exhibición, una de las sanciones previstas es el secuestro de la cosa, título o documento de que se trate.
\end{itemize}

\subsection{Secuestro previo a la subasta}

Una vez que el bien es embargado en el marco de un juicio de ejecución y es necesario proceder a la venta en pública subasta, por haberse agotado los pasos procesales previos, corresponde el secuestro como presupuesto de la subasta pública de los bienes embargados. Si el bien es un inmueble, no procede el secuestro, ya que este no recae directamente sobre inmuebles; pero los bienes muebles, registrables o no, tienen como paso previo a la subasta el secuestro.

\chapter{La intervención judicial}

\section{Concepto y fundamento}

\begin{definition}{Intervención judicial}
La intervención judicial (arts.\ 222 a 227 CPCCN) consiste en la designación, por parte del juez, de un auxiliar externo al tribunal para que se involucre en la administración o gestión de un ente colectivo, un patrimonio o una empresa. El interventor no tiene facultades de gobierno ni de disposición de bienes; su función es garantizar una administración adecuada o brindar información al juez.
\end{definition}

\section{Tipos de intervención judicial}

\subsection{Interventor recaudador}

El interventor recaudador (art. 223 CPCCN) es un auxiliar del tribunal designado por el juez, cuya misión es recaudar fondos del giro habitual de la administración de una empresa, un comercio o un patrimonio. Suele designarse cuando la parte deudora no tiene un patrimonio significativo, pero sí un comercio o consultorio: el interventor concurre cotidianamente al establecimiento y, del dinero que ingresa, retiene una suma para afectarla al pago de la deuda reclamada en el proceso.

\begin{example}{Empresa de transporte de ómnibus}
Una empresa de ómnibus de larga distancia que cobra sus ingresos por ventanilla (venta de pasajes) tiene una deuda de \$20.000. Embargar un ómnibus y sacarlo a subasta resulta un despropósito, porque el gasto de la subasta supera la deuda. Solución: se designa un interventor a ventanilla que retiene del ingreso cotidiano la suma correspondiente, logrando el cumplimiento de la deuda de forma proporcional y eficiente. En el caso de empresas de colectivos con tarjeta SUBE, puede comunicarse directamente por oficio a la administración de la gestión de las tarjetas para que retengan lo cobrado.
\end{example}

\subsection{Interventor informante y veedor}

El interventor informante (art. 224 CPCCN) está orientado a que el auxiliar que concurre al comercio de la empresa de la parte afectada dé noticia al juez acerca de la manera en que se lleva a cabo la administración, el patrimonio que posee la empresa, el estado de deudas que posee, y todo otro dato que el juez necesite. Esta información puede ser necesaria para decisiones que deban tomarse en el proceso.

\begin{example}{Sucesión con empresa}
En un juicio sucesorio, el patrimonio del causante incluye una empresa o comercio del que era titular. Para distribuir la herencia entre los herederos es necesario contar con información precisa sobre ingresos, egresos, bienes y deudas de ese patrimonio. En esos casos, la intervención del informante resulta necesaria para precisar el alcance de los bienes a distribuir y los ingresos que genera el patrimonio.
\end{example}

En doctrina se discute la diferencia entre el interventor informante y el veedor previsto en el art. 115 de la Ley 19.550 (Ley de Sociedades Comerciales), aunque la sutileza es fundamentalmente terminológica y no de alcance de gestión. En ambos casos, el alcance de la gestión y la información que debe brindarse queda establecido, más que por las normas del código, por lo que dispone el juez en la resolución que ordena la intervención.

\subsection{Interventor administrador o co-administrador}

En estos casos, el interventor gestiona ---en los casos puntuales que el juez comercial establece--- la administración, o cogestiona con administradores propios de la empresa esa administración. A diferencia del informante, aquí el interventor realiza cuestiones vinculadas con la administración cotidiana. Suele darse en casos en que la empresa se encuentra en estado prefallencial y es necesario evitar que, mediante maniobras de administración, se propicie la insolvencia a través del retiro de dinero por parte de los socios u otras maniobras perjudiciales.

\begin{important}
La intervención judicial no permite, en ninguna de sus formas, la disposición de bienes ni el gobierno de la empresa. Solo en el caso del administrador o co-administrador está habilitada la administración cotidiana. Los actos de disposición o de gobierno siempre deben ser adoptados por el juez, con intervención del síndico, cumpliendo los pasos procesales del proceso falencial. El administrador nunca tiene facultades para actos de disposición sin disposición expresa del juez.
\end{important}

\subsection{Quiénes pueden ser designados interventores}

Se aconseja que los interventores sean personas idóneas o técnicas en el negocio de que se trate. Generalmente se buscan entre administradores de empresa, contadores o incluso abogados para ciertas gestiones, como por ejemplo la recaudación en un comercio. Los abogados, contadores y licenciados en administración de empresas que pretenden ofrecer sus servicios para este tipo de intervenciones se anotan en los listados que elabora anualmente la Cámara Civil y la Cámara Comercial, presentando su currículum y poniéndose a disposición de los tribunales.

\chapter{La anotación de litis}

\section{Concepto y fundamento}

Este tema se conecta con el principio del tercero adquirente de buena fe a título oneroso estudiado en Derecho Civil. Cuando se declara la nulidad de un acto jurídico, uno de sus efectos es volver las cosas al estado anterior: ante la nulidad de un acto jurídico, todo vuelve al \textit{statu quo ante}. Si el acto jurídico era la compraventa o la donación de un inmueble y ese acto es declarado nulo, el inmueble vuelve al patrimonio de quien lo enajenó nulamente.

\begin{important}
\textbf{Art. 392 CCyCN --- tercero adquirente de buena fe a título oneroso}: cuando existe un tercero adquirente de buena fe a título oneroso, la ley privilegia su situación. Este tercero, que no fue parte en el acto declarado nulo, no tiene por qué asumir las consecuencias de operaciones viciadas anteriormente respecto del bien; por lo tanto, los efectos de la nulidad no se extienden a su respecto.
\end{important}

\begin{definition}{Anotación de litis}
Es una medida cautelar que, adoptada siempre sobre bienes registrables, publicita \textit{erga omnes} (es decir, con alcance a todos) la existencia de un proceso en trámite en el que se discute la validez o el alcance de una registración existente respecto de ese bien. Como consecuencia, al ser pública la existencia del juicio, nadie puede invocar luego la condición de adquirente de buena fe, porque nadie puede alegar buena fe cuando sabe que hay un litigio en trámite respecto del bien.
\end{definition}

Quien adquiere el bien en esas condiciones sabe que está adquiriendo un bien en riesgo de un litigio que puede resolverse en uno u otro sentido, y que, según cómo concluya ese juicio, tal vez deba devolverse a un anterior propietario que lo enajenó en virtud de un acto viciado.

\section{Requisitos de procedencia}

\begin{itemize}
    \item Solo procede contra bienes \textbf{registrables}.
    \item Debe existir una demanda iniciada. A diferencia del embargo, que puede trabarse antes del proceso, la anotación de litis exige en principio que exista juicio ya iniciado (se admite que el inicio de la mediación equivale al inicio del proceso).
    \item El juicio debe versar sobre una pretensión que, de prosperar, implicaría la modificación de un asiento registral respecto del bien (ya sea de titularidad u otra cuestión vinculada a lo que se anota registralmente).
\end{itemize}

\section{Supuestos típicos de procedencia}

\begin{table}[H]
\centering
\caption{Supuestos típicos de procedencia de la anotación de litis}
\begin{tabularx}{\textwidth}{l l X}
\toprule
\textbf{Tipo de juicio} & \textbf{Normativa} & \textbf{Particularidad} \\
\midrule
Usucapión / prescripción adquisitiva & Art. 1905 CCyCN & El juez ordena la anotación de oficio (sin pedido de parte) al disponer el traslado de la demanda. \\
\addlinespace
Expropiación & Arts. 24 y 14, Ley 21.499 & Una de las primeras medidas al iniciarse el proceso es la anotación de litis, para impedir que el propietario venda el bien y luego un tercero invoque buena fe frente a la expropiación. \\
\addlinespace
Escrituración & CPCCN & Busca modificar el asiento registral para que el adquirente por boleto pase a ser titular registral. \\
\addlinespace
Nulidad de escritura & Art. 392 CCyCN & Supuesto clásico de vicios del consentimiento o actos irregulares que intentan cancelar un asiento registral. \\
\addlinespace
Simulación & CCyCN & Caso típico de procedencia de la anotación de litis. \\
\addlinespace
Interdicto de adquirir & Art. 609 CPCCN & Hipótesis propia de la anotación de litis. \\
\bottomrule
\end{tabularx}
\end{table}

\section{Extinción de la anotación de litis}

La anotación de litis no se rige por el artículo 207 CPCCN. Se extingue por:

\begin{enumerate}
    \item \textbf{Cumplimiento de la sentencia}: si la demanda prospera, el cumplimiento de la sentencia determina la modificación registral buscada y concluye la anotación.
    \item \textbf{Desestimación del reclamo}: cuando el proceso concluye por sentencia desestimatoria, se agota la anotación de litis, debiendo comunicarse al registro.
    \item \textbf{Modos anormales de terminación}: caducidad de instancia, desistimiento o transacción. Cualquiera de ellos trae aparejada la extinción de la anotación de litis.
\end{enumerate}

\chapter{La prohibición de contratar}

\begin{note}
La prohibición de innovar (art. 230 CPCCN) se analiza en la Unidad 26 y será objeto de una explicación más detenida. En esta unidad se desarrolla directamente la prohibición de contratar (art. 231 CPCCN).
\end{note}

\section{Concepto y procedencia}

La prohibición de contratar (art. 231 CPCCN) es una cautela que, en determinados casos, resulta procedente, orientada a evitar que, con relación a un bien que es objeto del proceso o que se encuentra embargado, se realicen actos jurídicos que puedan disminuir o condicionar la cautela o el trámite del cumplimiento de la sentencia.

\begin{example}{Ejecución hipotecaria}
En el marco de un juicio de ejecución hipotecaria, además del embargo del bien, suele adoptarse la prohibición de contratar para evitar que el deudor hipotecario introduzca personas al inmueble mediante un contrato de locación o comodato, dificultando el avance de las medidas de cumplimiento de la sentencia en etapa de subasta. La prohibición de contratar puede notificarse también al locatario existente, para hacerle saber que no podrá prorrogarse ni ampliarse el contrato próximo a vencer.
\end{example}

\section{Notificación y publicidad}

La prohibición de contratar debe:

\begin{itemize}
    \item Si se trata de un bien registrable: asentarse en el registro para su publicidad \textit{erga omnes}.
    \item Notificarse al titular de la cosa o al propietario.
    \item Poder notificarse también a cualquier tercero que se sospeche pueda ser la persona que contrate con el propietario.
\end{itemize}

\begin{formula}{Caducidad especial de la prohibición de contratar (art. 231 CPCCN)}
La prohibición de contratar tiene una caducidad especial, más breve que la del art. 207 CPCCN: el plazo es de \textbf{cinco (5) días}. Este plazo corre desde el \textbf{dictado} de la medida (no desde su traba), y solo aplica si la prohibición se adopta \textbf{antes} del inicio del proceso. Si se adopta durante el curso del proceso ---incluso en la etapa conclusional--- no rige este plazo especial.
\end{formula}

\chapter{Protección de personas y cautelares en familia}

\section{La protección de personas y el artículo 234 CPCCN}

El análisis de esta medida obedece a una tradición que incluía una serie de potestades cautelares que poseía el juez en el antiguo artículo 234 CPCCN, referido a la protección de personas. Uno de sus incisos otorgaba amplias facultades al juez para tomar medidas de cambio de guarda y de disposición respecto de niños, en hipótesis en que se observara o se sospechara que estaban sufriendo maltrato, violencia o falta de cuidado.

\begin{important}
La \textbf{Ley 26.061} (Sistema Nacional de Protección Integral de Niñas, Niños y Adolescentes) derogó el inciso del art. 234 CPCCN que otorgaba facultades cautelares a los jueces respecto de niños. En consecuencia, sacó de la órbita del poder judicial toda la gestión de cuidado y protección de situaciones de maltrato, abuso o descuido de niñas, niños y adolescentes, y la trasladó al ámbito de la administración pública, a través de las Defensorías de Niños, Niñas y Adolescentes. Los llamados a intervenir ahora son la escuela, el hospital, los centros de salud, la policía y cualquier organismo que detecte una situación de riesgo.
\end{important}

Lo que sí procede y debe judicializarse, de acuerdo con los parámetros de esta ley, es el \textbf{control de legalidad} de las medidas que la administración (las defensorías) adopta respecto del cuidado o el abrigo de situaciones de riesgo de los niños. En tanto el órgano administrativo adopte medidas excepcionales de separación del niño de su ámbito familiar primario, estas medidas deben ser supervisadas y controladas en cuanto a su legalidad, para que no se prolonguen en el tiempo más allá de lo necesario.

\begin{note}
El grueso de la actividad protectora que poseía anteriormente el art. 234 CPCCN ha quedado transferido a la administración pública, a través del Sistema Nacional de Protección Integral de Niños, Niñas y Adolescentes (Ley 26.061). El artículo solo conserva una función tuitiva de los jueces en materia de mayores de edad con restricciones a la capacidad por problemáticas de salud mental.
\end{note}

\section{Potestades cautelares en materia de violencia}

\begin{table}[H]
\centering
\caption{Regímenes cautelares en materia de violencia familiar y de género}
\begin{tabularx}{\textwidth}{l X}
\toprule
\textbf{Norma} & \textbf{Contenido} \\
\midrule
Ley 24.417 --- Ley de Violencia Familiar & Establece medidas cautelares que los jueces pueden adoptar para erradicar cualquier tipo de situación de agresión padecida por personas mayores de edad (hombres, mujeres y niños involucrados con personas mayores de edad). Incluye medidas de distanciamiento, prohibición de acercamiento, fijación provisional de alimentos, regímenes de comunicación y cuidado personal de los niños. \\
\addlinespace
Ley 26.485 --- Ley de Protección contra la Violencia hacia las Mujeres & Refuerza y amplía las medidas de la Ley 24.417 en materia de perspectiva de género. Otorga al sistema judicial potestades cautelares de amplísima gama: distanciamiento, prohibición de acercamiento, prestación alimentaria, fijación de regímenes de comunicación entre padres e hijos, establecimiento del cuidado personal de los niños, y medidas patrimoniales vinculadas a la distribución de bienes. \\
\bottomrule
\end{tabularx}
\end{table}

\section{Medidas cautelares en el proceso de familia}

En materia de derecho de familia, la vara que evalúa los presupuestos de verosimilitud del derecho, peligro en la demora y contracautela está superada por el hecho del vínculo familiar existente entre las personas involucradas.

\begin{table}[H]
\centering
\caption{Aplicación de los presupuestos cautelares en el proceso de familia}
\begin{tabularx}{\textwidth}{l X}
\toprule
\textbf{Presupuesto} & \textbf{Aplicación en procesos de familia} \\
\midrule
Verosimilitud del derecho & Surge del propio vínculo familiar (conyugal, de pareja, filial). No es necesario evaluar más que la existencia del vínculo como presupuesto que satisface este recaudo. \\
\addlinespace
Peligro en la demora & Está dado \textit{in re ipsa} (en la propia naturaleza de las cosas). La urgencia de cautelar el cuidado de los hijos, los alimentos provisorios, la vivienda y el régimen de comunicación es evidente por los derechos fundamentales comprometidos. \\
\addlinespace
Contracautela & En la generalidad de los casos no existe contracautela exigible. La caución juratoria está implícita y ni siquiera se pide, ya que resulta innecesaria. \\
\bottomrule
\end{tabularx}
\end{table}

\subsection{Las cautelares de familia en el Código Civil y Comercial de la Nación}

Las medidas cautelares de familia se encuentran, en su gran mayoría, insertas en el código de fondo, el Código Civil y Comercial de la Nación. Esto no es una novedad introducida por la reforma de 2015, sino que ya venía con el código anterior. El legislador siempre ha considerado que ciertos aspectos del conflicto familiar competen al legislador nacional garantizar, sin dejarlos librados a lo que pueda establecer el legislador local provincial. Por ello plasmó en el código de fondo normas operativas vinculadas con medidas cautelares propias del derecho de familia, que abarcan:

\begin{itemize}
    \item La atribución del hogar familiar.
    \item La garantía de subsistencia alimentaria provisoria mientras se define, en un proceso de conocimiento, las prestaciones alimentarias para los hijos menores y, eventualmente, entre cónyuges.
    \item La definición de la manera en que los padres ejercerán el cuidado personal de los hijos.
    \item El régimen de comunicación de los hijos con sus padres.
    \item Las medidas de carácter patrimonial vinculadas con la distribución de bienes como consecuencia de la extinción del régimen matrimonial.
\end{itemize}

\begin{note}
Como lectura recomendada se sugiere el artículo publicado en el sitio web www.jorgearojas.com.ar: ``Las medidas cautelares en el Código Civil y Comercial de la Nación'', del Dr.\ Jorge A.\ Rojas, que analiza las medidas cautelares nominadas e innominadas del CCyCN, las heredadas del código anterior, las absolutamente novedosas, y aspectos de tutelas anticipatorias contempladas en el código civil y comercial.
\end{note}

% ============================================================
% CORNELL — REPASO GENERAL DE LA UNIDAD
% ============================================================
\chapter{Cuadro Cornell: repaso general de la unidad}

El siguiente cuadro reúne, en formato Cornell, las preguntas clave que permiten repasar de manera activa los contenidos centrales de la unidad, junto con su desarrollo y una síntesis final.

\begin{longtable}{p{0.30\textwidth} p{0.62\textwidth}}
\toprule
\textbf{Preguntas / palabras clave} & \textbf{Notas / respuestas} \\
\midrule
\endfirsthead

\toprule
\textbf{Preguntas / palabras clave} & \textbf{Notas / respuestas} \\
\midrule
\endhead

\midrule
\multicolumn{2}{r}{\textit{Continúa en la página siguiente}} \\
\endfoot

\bottomrule
\endlastfoot

\textbf{¿Qué es el embargo y por qué se lo llama ``medida cautelar por antonomasia''?} &
Es la medida cautelar más usada y habitual. Consiste en la afectación e individualización de un bien (mueble o inmueble, registrable o no) al resultado de un proceso. El bien queda vinculado al proceso y opera como garantía concreta de la efectividad de la sentencia. Se lo denomina ``por antonomasia'' porque es la más frecuente en la práctica judicial. \\
\addlinespace

\textbf{¿Cuáles son los tres tipos de embargo?} &
1) Embargo \textbf{preventivo}: eminentemente cautelar, puede trabarse antes del proceso o durante su tramitación. Exige los tres recaudos clásicos (verosimilitud, peligro en la demora, contracautela). Arts. 209-212 CPCCN. 2) Embargo \textbf{ejecutivo}: se traba una vez que el título ha sido declarado admisible en el juicio ejecutivo y se libra la intimación de pago. 3) Embargo \textbf{ejecutorio}: se traba en el proceso de ejecución de sentencia; es el paso previo a la subasta y opera automáticamente por ministerio de la ley (art. 551 CPCCN). \\
\addlinespace

\textbf{¿Puede venderse un bien embargado? ¿Cuáles son las consecuencias?} &
Sí, puede venderse: el embargo no impide la enajenación. Pero por el principio \textit{nemo plus iuris}, nadie puede transmitir un derecho mayor al que posee: quien vende transfiere el bien con el embargo, y el adquirente asume el embargo. Según el Plenario ``Sertop'' (Cámara Nacional Civil), el adquirente responde por la totalidad de la deuda ---capital, ampliaciones, intereses y costas--- no solo por el monto de la traba. En el fuero comercial este criterio no rige. \\
\addlinespace

\textbf{¿Cuáles son los principales bienes inembargables?} &
Art. 219 CPCCN: lecho cotidiano, ropa y bienes de uso indispensable, sepulcros, y bienes declarados inembargables por leyes especiales. La Ley de Contrato de Trabajo limita el embargo del salario al 20\% del ingreso (con excepción en deudas alimentarias). Las jubilaciones y pensiones también son inembargables en principio, salvo deudas alimentarias. \\
\addlinespace

\textbf{¿Cuándo caduca el embargo?} &
Si se trabó antes del proceso, caduca si no se promueve el proceso principal dentro de los 10 días (art. 207 CPCCN). El inicio de la mediación interrumpe el plazo. Los embargos registrales en el Registro del Automotor caducan a los cinco años y deben reinscribirse para mantenerse vigentes. \\
\addlinespace

\textbf{¿Qué es el secuestro y sobre qué bienes recae?} &
El secuestro es una medida cautelar que desapodera al deudor de la disposición material y jurídica sobre una cosa. Solo opera sobre bienes \textbf{muebles} (registrables o no). Tiene tres modalidades: (1) cautelar propiamente dicha (art. 221 CPCCN), (2) como sanción por incumplimiento (arts. 128 y 323 CPCCN), y (3) como paso previo a la subasta de bienes muebles embargados. \\
\addlinespace

\textbf{¿Qué facultades tiene el interventor judicial?} &
El interventor no tiene facultades de gobierno ni de disposición de bienes. Hay tres tipos principales: (1) \textbf{Recaudador} (art. 223): retiene fondos del giro normal de la empresa para pagar la deuda. (2) \textbf{Informante} (art. 224): aporta información al juez sobre la administración y el patrimonio. (3) \textbf{Administrador / co-administrador} (Ley 19.550): solo en casos de intervención comercial especial, gestiona la administración cotidiana sin poder disponer bienes. \\
\addlinespace

\textbf{¿Qué es la anotación de litis y qué protege?} &
Es una medida cautelar que publicita \textit{erga omnes} ---sobre bienes registrables--- la existencia de un proceso en trámite donde se discute la validez de un asiento registral. Su efecto principal es impedir que alguien invoque la buena fe del art. 392 CCyCN como tercero adquirente, porque quien compra sabiendo del litigio no puede alegar ignorancia. Solo procede sobre bienes registrables y exige juicio iniciado. \\
\addlinespace

\textbf{¿En qué supuestos procede la anotación de litis de oficio?} &
En los juicios de usucapión (art. 1905 CCyCN), el juez ordena la anotación de oficio al disponer el traslado de la demanda, sin necesidad de pedido de parte. También opera de oficio en los juicios de expropiación (art. 24, Ley 21.499). \\
\addlinespace

\textbf{¿Cómo se extingue la anotación de litis?} &
Se extingue: (1) con el cumplimiento de la sentencia favorable al actor (se produce la modificación registral buscada), (2) con la sentencia desestimatoria, y (3) por modos anormales de terminación del proceso (caducidad, desistimiento, transacción). En todos los casos debe comunicarse la extinción al registro. \\
\addlinespace

\textbf{¿Para qué sirve la prohibición de contratar?} &
Está orientada a impedir que, respecto de un bien objeto del proceso o embargado, se realicen actos jurídicos que puedan disminuir o condicionar la cautela o el cumplimiento de la sentencia. Se usa frecuentemente en ejecuciones hipotecarias para evitar que el deudor introduzca locatarios que dificulten la subasta. Tiene una caducidad especial de 5 días si se adopta antes del inicio del proceso. \\
\addlinespace

\textbf{¿Qué cambió respecto de la protección de niños con la Ley 26.061?} &
La Ley 26.061 derogó el inciso del art. 234 CPCCN que otorgaba a los jueces amplias facultades para dictar medidas de cambio de guarda y abrigo de niños. Trasladó esa competencia a la administración pública (Defensorías de NNA). Los jueces ahora solo ejercen el control de legalidad de las medidas excepcionales que adopta la administración. El art. 234 subsiste solo para mayores con restricción de capacidad por salud mental. \\
\addlinespace

\textbf{¿Qué particularidades tienen los presupuestos de las cautelares en familia?} &
Los tres presupuestos clásicos (verosimilitud, peligro en la demora, contracautela) están prácticamente dados \textit{in re ipsa}: (1) la verosimilitud surge del propio vínculo familiar; (2) el peligro en la demora es inherente a la naturaleza urgente de las necesidades familiares (alimentos, cuidado, vivienda); (3) la contracautela es prácticamente inexistente, ya que la caución juratoria está implícita y ni siquiera se exige. \\
\addlinespace

\multicolumn{2}{p{0.94\textwidth}}{
\textbf{Resumen:} La unidad desarrolla las medidas cautelares en particular del proceso civil, todas articuladas sobre los tres presupuestos generales: verosimilitud del derecho, peligro en la demora y contracautela. El embargo es la medida cautelar por antonomasia ---la más habitual--- y su correcta comprensión exige distinguir sus tres modalidades según el estadio procesal en que se traba (preventivo, ejecutivo, ejecutorio), conocer sus efectos reales sobre la circulación de bienes (principio \textit{nemo plus iuris}, Plenario Sertop) y sus límites de embargabilidad (bienes inembargables, límite del 20\% del salario). El secuestro, más excepcional, despoja al deudor de la tenencia material del bien mueble y asume tres formas según su finalidad (cautelar, sanción o presupuesto de subasta). La intervención judicial coloca a un auxiliar del tribunal en la gestión de una empresa o patrimonio ---sin facultades de disposición--- en sus variantes de recaudador, informante/veedor y administrador. La anotación de litis opera sobre bienes registrables como mecanismo de publicidad que neutraliza la buena fe del tercero adquirente, siendo exigida de oficio en juicios de usucapión y expropiación. La prohibición de contratar protege la integridad del bien durante el proceso, con una caducidad especial de cinco días si es previa al inicio. Finalmente, en el ámbito de familia, la Ley 26.061 trasladó la adopción de medidas de protección de niños a la administración pública ---reservando a los jueces solo el control de legalidad--- mientras que las leyes 24.417 y 26.485 otorgan al Poder Judicial amplísimas potestades cautelares en materia de violencia familiar y de género. En los procesos de familia, los presupuestos cautelares se tienen por configurados \textit{in re ipsa}, derivados de la propia naturaleza urgente del conflicto familiar.
} \\

\end{longtable}

\end{document}
